\chapter*{Introduction}
\addcontentsline{toc}{chapter}{Introduction}
\markboth{Introduction}{Introduction}

The transition towards Post-Quantum Cryptography (PQC) is driven by the potential of quantum computers to break traditional asymmetric encryption.
Standard algorithms like RSA and ECC are vulnerable to Shor's algorithm~\cite{shor1994algorithms}, which solves the underlying mathematical problems in polynomial time.
In response, the National Institute of Standards and Technology (NIST) has spearheaded a process to standardize algorithms resistant to quantum attacks~\cite{fips203}.

Lattice-based cryptography is a primary focus of this standardization, particularly schemes based on the Learning With Errors (LWE) problem.
While LWE offers high security, its implementation in practical systems is often hindered by large memory requirements $O(n^2 \log q)$ for keys, contributing to kilobyte-sized public keys for modest security levels~\cite{Lindner2011-wg}.
The Module-LWE (M-LWE) variant addresses these limitations by utilizing structures (such as polynomial rings) that allow for smaller keys and faster computations while maintaining similar security levels to regular LWE\@.

This thesis explores the optimization and efficient software implementation of Module-LWE schemes. The focus is placed on the mathematical structures used in the CRYSTALS-Kyber algorithm~\cite{kyber_spec_2021}, analyzing how the Number Theoretic Transform (NTT) and Single Instruction, Multiple Data (SIMD) vector instructions can be leveraged to achieve high-speed execution. By optimizing these core components, this work demonstrates that Module-LWE provides a robust and efficient solution for secure digital communication in the post-quantum era.

The primary objective is to implement and optimize the core operations of Module-LWE based cryptographic schemes, specifically targeting the Kyber-512 parameter profile. This involves analyzing the mathematical foundations of the cryptosystem, developing an optimized software implementation in C, and conducting a detailed performance evaluation on modern hardware.

The structure of this thesis is as follows. \textbf{Chapter \ref{ch:theoretical_background} (Theoretical Background)} provides the mathematical foundations of lattice-based cryptography, the quantum threat, and the transition from standard LWE to Module-LWE. \textbf{Chapter \ref{ch:chosen_methods} (Chosen Methods)} details the algorithmic optimizations, specifically focusing on the Number Theoretic Transform and the requirements for SIMD-accelerated arithmetic. \textbf{Chapter \ref{ch:implementation} (Implementation)} describes the software architecture of the developed C11 library and the implementation of the benchmarking and verification suites. \textbf{Chapter \ref{ch:performance_evaluation} (Performance Evaluation)} presents the results of benchmarks conducted on modern hardware, comparing the optimized implementation against reference standards. Finally, the concluding chapter summarizes the empirical findings and discusses potential paths for future research in post-quantum cryptographic software optimization.
